\minitopic{专题研讨：证成的时间结构与击败机制}证成理论常被画成静止图：底层经验支持中层信念，中层信念再支持结论。但真实探究具有时间。主体先因证据形成信念，后来忘记部分理由，遇到质疑，查询记录，或发现自己的判断过程受偏见影响。一个完整理论要解释信念怎样获得证成、怎样保持证成、何时失去证成以及如何恢复。只问“此刻有哪些理由支持命题”会漏掉信念所依据的路径；只问过程过去是否可靠，又会漏掉当前已经可得的反证。

\minitopic{从支持关系到基于关系}命题 \(p\) 可能被主体拥有的证据 \(e\) 充分支持，但主体相信 \(p\) 的原因却是愿望、猜测或权威压力。此时存在命题确证而缺少信念确证。所谓基于关系，旨在连接“有一个好理由”与“因为这个理由而相信”。单纯因果说要求 \(e\) 实际造成信念；反事实说要求若 \(e\) 改变，信念相应改变；理由回应说则要求主体能以适当方式调整信念，而不一定能口头重建全部推理。 三种说法各有反例。药物让受试者更专注，从而注意到实验数据；数据造成信念，药物也是因果条件，但只有数据是其认识理由。学生背出老师的论证，却无论前提怎样都坚持答案，说明表面引用不等于反事实依赖。专家一眼识别病灶，可能无法立刻说明全部线索，却仍由长期训练形成的模式识别适当回应证据。基于关系因而既不能缩成任意原因，也不能一律要求可陈述推理。 研究时可以构造“交换理由”测试：保持结论不变，把原证据替换为同样支持结论但内容不同的证据，再问主体的信心、解释和行动是否改变。也可构造“移除支持”测试：若删除被声称为理由的材料，信念仍毫无变化，理由可能只是事后合理化。但这些测试给出的是诊断信号，并非充分条件；坚定信念有时由多个冗余理由支持。

\minitopic{回溯、无穷与循环}认识回溯提出三难：每个理由都需要新理由，似乎只能无限延伸、在某处停止或形成循环。强基础主义让不可错经验终止回溯，却难以说明普通知觉为何足够；温和基础主义允许经验提供可错的初步支持，同时接受击败。融贯主义把支持放在系统的相互联系中，但必须解释系统为何不是一篇内部一致的小说。\parencite{bonjour1985} “停止”并不自动等于独断。理由追问总在某个问题背景中发生。看见仪器读数可以初步支持温度判断；若有人提供校准失效证据，就需要进一步检查。基础性可以被理解为默认资格，而非绝对免疫。反之，“彼此支持”也不必是狭窄的两项循环。观察、模型、独立测量和预测成功构成广泛网络时，新增节点可能提高整体解释力；不过若所有输入都来自同一故障源，网络规模不会创造独立性。 一个实用的论证图应标出边的类型：演绎支持、概率支持、解释支持、来源认证和击败解除。把所有箭头都画成“支持”会掩盖循环。例如，仪器可靠性由标准样本检验，标准样本的性质又完全由该仪器确定，这是恶性循环；若标准样本还有不同原理的测量和制造记录，循环可能被外部约束。结构判断离不开来源独立性。

\minitopic{反驳型与削弱型击败}反驳型击败者提供相信非 \(p\) 的理由；削弱型击败者不支持非 \(p\)，而是破坏原证据与 \(p\) 之间的联系。看到天气预报改为晴天，是“明天下雨”的反驳；得知预报网站缓存了上周页面，则削弱原预报，但本身不说明明天天气。混淆二者会导致过度更新：来源失效后，合理态度往往是暂停，而不是立刻相信相反命题。 击败还分为实际存在、主体可得与主体应当取得。档案库里有一封反证邮件，但研究者没有合理机会发现；同事当面指出编码错误，研究者却故意不读；匿名账户只说“你的结论错了”，没有给出可查理由。内在主义倾向强调主体可把握的理由，外在主义更关注过程与事实的适当联系。\parencite{goldman1979,conee1985} 规范评价可以分层：信念是否知识导向、主体是否主观合理、是否尽到查询责任，三者未必同值。 解除击败不是把不利信息遗忘。若同事质疑样本标签，研究者用独立日志确认无误，原击败被回应；若只发现同事过去常犯错，这可能降低其可信度，但仍要说明本次具体警报为何可以忽略。解除理由必须针对击败机制。压低消息来源的社会地位、切断提醒频道或反复自我保证，都不能恢复原证据的资格。

\minitopic{高阶证据与认识自信}一阶证据直接涉及世界是否为 \(p\)，高阶证据涉及自己是否正确评估一阶证据。你完成一份概率计算，随后得知自己服用了会显著损害运算能力的药物；药物信息不直接说明答案是什么，却挑战你对计算的信赖。遇到同等能力者的分歧也可能产生高阶证据。此时“我仍看不出错误”不足以解除挑战，因为受损能力恰好能解释为何看不出。 严格调和论要求大幅降低信心，坚定论允许在某些条件下维持原判断。争点不宜抽象处理：双方是否真正共享证据，能力是否在本题相关，是否存在独立核验，分歧能否由不同先验或价值权衡解释？许可主义进一步认为同一证据有时允许多个理性态度。许可范围过宽会让证据失去约束；范围过窄则假定每个复杂问题都有唯一精确置信度。可操作的中间立场要求态度落在受证据限制的区间，并对发现的新分歧作出二阶更新。 高阶反思也可能陷入“为什么相信自己的方法可靠”的循环。主体不能在每次推理前完成对全部认知能力的独立证明。合理要求不是先验地认证整个心智，而是在出现具体故障迹象时使用相对独立的校验：换一种算法、让同事盲审、检查原始数据或延迟决定。反思的价值在于发现局部脆弱性，不在于获得不可错的自我许可证。

\minitopic{比较视野：学、思与可改正性}《论语》把学与思并置，可以被用来提出一个结构问题：外来材料与主体整理何以相互约束？若只有接受而无辨析，信念可能缺少基于关系；若只有内部思索而无对象约束，系统可能融贯却悬空。\parencite{analects2003} 这与当代基础主义或融贯主义存在功能上的交叉，但不意味着古典文本已经提出同一组技术定义。比较的价值是提醒我们，证成也包含学习实践中的可教、可问和可改正。 荀子强调通过师法、积习和辨蔽改造认识能力。\parencite{xunzi2014} 从外在主义看，制度训练可以提高过程可靠性；从德性认识论看，它塑造注意、耐心与回应批评的能力；从内在主义看，主体仍需逐步理解规则为何构成理由。三种解读不必相互排斥，却评价不同层面。比较研究应明确：文本是在说明信念为真、主体有责、能力形成，还是共同体可传承知识。

\minitopic{综合案例：被撤回的数据集}一项公共卫生研究依据开放数据集得出药物有效结论。分析代码正确，同行也复现了结果。六个月后，数据机构公告：部分医院上传时把“接受治疗”和“对照”标签对调；公告没有说明受影响医院名单。研究者此时应如何评价原信念？ 首先区分时间：最初根据当时可得信息形成结论，可能主观合理；公告出现后，原信念受到削弱型击败，因为数据—结论联系不再稳固，但尚未获得药物无效的证据。其次区分对象：代码可靠性没有被击败，数据语义被击败，复现若使用同一数据并非独立确认。再次区分责任：研究者应暂停强断言、请求医院级审计并测试最坏情形，而不是删除旧论文或假装已知相反结论。 若后来审计显示对调只发生在未纳入样本的医院，击败得到针对性解除；若研究者只以“结果已通过同行评议”为由维持信心，同行评议并未触及新故障机制。请用论证图标出一阶支持、来源认证、高阶警报与解除路径，再分别从内在主义、可靠主义和责任评价给出判断。一个好答案可以承认：同一主体在某时刻“没有应受责备”，却“已经不再知道”。

\begin{tcolorbox}[breakable,colback=white,colframe=blue!30!black,title={专题研读任务},fonttitle=\bfseries]
选取一项近期被更正的研究或统计结论，建立时间轴：信念形成、首次反证、调查、解除或撤回。对每个节点分别记录命题确证、信念确证、知识、主观合理性与查询义务，避免用一个“可信／不可信”标签覆盖全过程。
\end{tcolorbox}

\index{基于关系}
\index{认识回溯}
\index{击败者}
\index{高阶证据}
\index{许可主义}
